
% =========================================================
% MANUAL RÁPIDO DE USO DO BIBTEX NO MODELO DE TESE UFPI
% =========================================================
%
%
% ---------------------------------------------------------
% 1. ARQUIVOS ENVOLVIDOS
% ---------------------------------------------------------
%
% Neste modelo, a bibliografia é controlada principalmente
% por dois arquivos:
%
% 1) referencias_tese.bib
%    -> guarda as entradas bibliográficas.
%
% 2) tese_ufpi_revisada.tex
%    -> arquivo principal, que chama a bibliografia no final.
%
% No arquivo principal, costumam aparecer linhas como:
%
% \bibliographystyle{plain}
% \bibliography{referencias_tese}
%
% A primeira linha escolhe o estilo da bibliografia.
% A segunda informa o nome do arquivo .bib.
%
% ---------------------------------------------------------
% 2. COMO INSERIR UMA REFERÊNCIA NO ARQUIVO .BIB
% ---------------------------------------------------------
%
% Cada item bibliográfico recebe uma chave única.
% Exemplo:
%
% @article{grigoryan1999,
%   author  = {Grigor'yan, Alexander},
%   title   = {Analytic and geometric background of recurrence and non-explosion of the Brownian motion on Riemannian manifolds},
%   journal = {Bulletin of the American Mathematical Society},
%   volume  = {36},
%   number  = {2},
%   pages   = {135--249},
%   year    = {1999}
% }
%
% Aqui, a chave é: grigoryan1999
%
% É essa chave que será usada em:
%
% \cite{grigoryan1999}
%
% ---------------------------------------------------------
% 3. COMO CITAR NO TEXTO
% ---------------------------------------------------------
%
% Para citar uma referência no corpo da tese, use por exemplo:
%
% Como em \cite{grigoryan1999}, ...
%
% Para citar mais de uma referência ao mesmo tempo:
%
% Ver \cite{grigoryan1999,pessoa}.
%
% ---------------------------------------------------------
% 4. TIPOS MAIS COMUNS DE ENTRADA
% ---------------------------------------------------------
%
% Os formatos mais usados no arquivo .bib são:
%
% @article
% -> artigo publicado em periódico
%
% @book
% -> livro
%
% @incollection
% -> capítulo de livro
%
% @misc
% -> preprint, nota técnica, material sem formato editorial fixo
%
% Exemplo de livro:
%
% @book{evans,
%   author    = {Evans, Lawrence C.},
%   title     = {Partial Differential Equations},
%   publisher = {American Mathematical Society},
%   year      = {2010}
% }
%
% Exemplo de preprint:
%
% @misc{pessoa,
%   author        = {Pessoa, Leandro F. and Pigola, Stefano and Setti, Alberto G.},
%   title         = {Dirichlet parabolicity and {$L^{1}$}-Liouville property under localized geometric conditions},
%   howpublished  = {arXiv preprint},
%   eprint        = {1607.06483},
%   archivePrefix = {arXiv},
%   primaryClass  = {math.DG},
%   year          = {2017}
% }
%
% ---------------------------------------------------------
% 5. ORDEM DE COMPILAÇÃO
% ---------------------------------------------------------
%
% Quando o documento usa BibTeX, a sequência típica é:
%
% pdflatex tese_ufpi_revisada
% bibtex tese_ufpi_revisada
% pdflatex tese_ufpi_revisada
% pdflatex tese_ufpi_revisada
%
% No VS Code, isso pode ser automatizado pela extensão
% LaTeX Workshop, desde que a receita de compilação inclua BibTeX.
%
% ---------------------------------------------------------
% 6. CUIDADOS IMPORTANTES
% ---------------------------------------------------------
%
% - Não repetir a mesma chave em duas entradas diferentes.
% - Usar uma chave curta e fácil de reconhecer.
% - Preencher os campos essenciais: autor, título e ano.
% - Em artigos, incluir também periódico, volume, número e páginas.
% - Em preprints, incluir eprint, archivePrefix e primaryClass.
%
% ---------------------------------------------------------
% 7. SUGESTÃO DE PADRÃO PARA AS CHAVES
% ---------------------------------------------------------
%
% Um padrão simples e limpo:
%
% sobrenomeAno
%
% ou, quando houver vários autores:
%
% primeiroAutorAno
%
% Exemplos:
% grigoryan1999
% pigola2013
% holopainen1990
%
% ---------------------------------------------------------
% 8. EXEMPLO COMPLETO DE USO NO TEXTO
% ---------------------------------------------------------
%
% Suponha que o arquivo referencias_tese.bib contenha a entrada:
%
% grigoryan1999
%
% Então, no capítulo, você pode escrever:
%
% A discussão analítica de recorrência e não explosão pode ser encontrada em \cite{grigoryan1999}.
%
% Ao compilar corretamente, a citação aparecerá no texto
% e a referência completa será listada ao final da tese.
%
% ---------------------------------------------------------
% OBSERVAÇÃO FINAL
% ---------------------------------------------------------
%
% Este arquivo foi pensado como manual de leitura direta no código.
% Portanto, o conteúdo principal está em comentários.
% Se quiser, depois você pode manter este arquivo apenas como
% documentação interna do modelo.
%
